\centerline{\Huge \bf  Подобие на стероидах}
\centerline{\Huge \bf Гомотетия помогает решать задачи}

\begin{flushright}
    \scriptsize{}
\end{flushright}


\textit{Определение}. Гомотетией с центром в точке $O$ и коэффициентом $k \neq 0$ называется преобразование плоскости, при котором точка $A$ переходит в точку $A'$, которая лежит на прямой $OA$ и $OA'$ = $k \cdot OA$, причём при $k > 0$ точки $A$ и $A'$ лежат по одну, а при $k < 0 \ -$ по разные стороны от точки $O.$

\noindent\textbf{Свойства гомотетии.} Докажите, что гомотетия:

а) переводит прямые в прямые, окружности в окружности; 

б) сохраняет отношения длин отрезков и углы между прямыми; 

в) сохраняет длины отрезков (является движением) тогда и только тогда, когда её коэффициент равен $1$ или $-1$.

\problem{1}
Докажите, что композиция (результат последовательного применения) двух гомотетий с коэффициентами $k_1$ и $k_2$, где $k_1 \cdot k_2 \neq 1$, есть гомотетия, причём центры всех трёх гомотетий лежат на одной прямой. Что будет в случае $k_1 \cdot k_2 = 1$?

\problem{2}
Докажите, что в любой трапеции середины оснований, точка пересечения диагоналей и точка пересечения боковых сторон лежат на одной прямой.

\problem{3}
Две окружности касаются в точке $K$. Прямая, проходящая через точку $K$, пересекает эти окружности в точках $A$ и $B$. Докажите, что касательные к окружностям, проведенные через точки $A$ и $B$, параллельны.

\problem{4}
Две окружности касаются в точке $K$. Через точку $K$ проведены две прямые, пересекающие первую окружность в точках $A$ и $B$, вторую $-$ в точках $C$ и $D$. Докажите, что $AB \parallel CD$.

\problem{5} На сторонах $AB$ и $AC$ треугольника $ABC$ взяты соответственно точки $M$ и $N$, причём  $MN \parallel BC$.  На отрезке MN взята точка $P$, причём  $MP = \dfrac{1}{3}MN.$  Прямая $AP$ пересекает сторону $BC$ в точке $Q$. Докажите, что  $BQ = \dfrac{1}{3}BC$.